\section*{Hoofdstuk \ref{ch:g9:maps-of-the-earth} --- Kaarten van de hele Aarde}

\begin{solution}{exo:g9:maps-of-the-earth:1}
(a) Een \omterm{def:g9:maps-of-the-earth:projection}{kaartprojectie}. (b) Een \omterm{def:g9:maps-of-the-earth:contour}{hoogtelijn}. (c) Een \omterm{def:g9:maps-of-the-earth:numerical-scale}{getalschaal}. (d) De
\omterm{def:g9:maps-of-the-earth:contour}{equidistantie}. (e) \omterm{def:g9:maps-of-the-earth:remote-sensing}{Teledetectie}.
\end{solution}

\begin{solution}{exo:g9:maps-of-the-earth:2}
(a) $4 \times 25\,000 = \num{100000}$~cm $= \qty{1}{km}$. (b) $3 \times
50\,000 = \num{150000}$~cm $= \qty{1.5}{km}$. (c) $2.5 \times 1\,000\,000
= \num{2500000}$~cm $= \qty{25}{km}$. (d) $\qty{1}{km} = \num{100000}$~cm;
$100\,000 \div 25\,000 = \qty{4}{cm}$.
\end{solution}

\begin{solution}{exo:g9:maps-of-the-earth:3}
(a) Niet waar: een bol kan niet plat worden gemaakt zonder lengtes te
vervormen, dus elke \omterm{def:g3:continents-and-oceans:world-map}{wereldkaart} vervormt iets. (b) Waar. (c) Niet waar: een
oppervlaktegetrouwe \omterm{def:g2:drawing-to-map:map}{kaart} behoudt oppervlakten, maar vervormt de vormen,
vooral bij haar randen. (d) Waar.
\end{solution}

\begin{solution}{exo:g9:maps-of-the-earth:4}
(a) De helling is steil. (b) Het terrein loopt flauw af. (c) Een \omterm{def:g1:landscapes-around-us:landform}{heuvel},
met de top binnen de kleinste lijn. (d) Een punt van het terrein heeft maar
één hoogte, dus het kan niet op twee lijnen met verschillende hoogten
liggen.
\end{solution}

\begin{solution}{exo:g9:maps-of-the-earth:5}
(a) Op \ang{60}\,NB wordt het \omterm{def:g1:landscapes-around-us:water}{meer} 4 keer te groot getekend: het lijkt zo
groot als $4 \times 1000 = \qty{4000}{km^2}$. Op de \omterm{def:g3:continents-and-oceans:equator}{evenaar} worden
oppervlakten niet vergroot, dus het \omterm{def:g1:landscapes-around-us:water}{meer} dat even groot is getekend, heeft
een ware oppervlakte van \qty{4000}{km^2}. (b) $4 \div 1.33 \approx 3$ keer
groter. (c) Oppervlakten, en afstanden van noord naar zuid, groeien naar de
\omterm{def:g3:continents-and-oceans:poles}{polen} toe onbeperkt: de \omterm{def:g3:continents-and-oceans:poles}{polen} zouden oneindig ver van de \omterm{def:g3:continents-and-oceans:equator}{evenaar} liggen.
\end{solution}

\begin{solution}{exo:g9:maps-of-the-earth:6}
(a) \qty{50}{m}. (b) De top ligt op \qty{544}{m}; de hoofdhoogtelijn
eromheen is de lijn van \qty{500}{m}. (c) In het westen: een \omterm{def:g1:landscapes-around-us:water}{rivier} stroomt
bergafwaarts, dus haar bodem ligt stroomafwaarts, naar het oosten, lager.
(d) Tussen ongeveer de lijnen van 300 en \qty{400}{m}: daar is de helling
het steilst, en ze is flauwer zowel bij de dalbodem als bij de top.
\end{solution}

\begin{solution}{exo:g9:maps-of-the-earth:7}
(a) C op ongeveer \qty{400}{m}, D op ongeveer \qty{530}{m}. (b) Ongeveer
\qty{135}{m}, ongeveer \qty{1.3}{km} van C, op de dalbodem. (c) Ongeveer
0.3 tot \qty{0.4}{km}. (d) De hoogteverschillen zijn klein vergeleken met
de afstanden: zonder uitrekken zouden de \omterm{def:g1:landscapes-around-us:landform}{heuvels} bijna vlak lijken.
\end{solution}

\begin{solution}{exo:g9:maps-of-the-earth:8}
(a) $2 \times 25\,000 = \num{50000}$~cm $= \qty{500}{m}$. (b) $300 - 200 =
\qty{100}{m}$. (c) $100 \div 500 = 0.2$, een helling van \qty{20}{\%}. (d) Minder steil:
\qty{4}{cm} is \qty{1000}{m}, dus de helling is $100 \div 1000 = \qty{10}{\%}$,
half zo steil.
\end{solution}

\begin{solution}{exo:g9:maps-of-the-earth:9}
(a) $14 \times 18 = 252$ omlopen. (b) $14 \times 103 = 1442$ minuten,
ongeveer \qty{24}{h}: een dag. (c) $5.5 \div 1 = 5.5$ keer. (d) Bijvoorbeeld:
een \omterm{def:g8:ice-on-earth:glacier}{gletsjer} die zich terugtrekt, een bos dat wordt gekapt, een stad die
groeit, een \omterm{def:g1:landscapes-around-us:water}{kust} of een \omterm{def:g1:landscapes-around-us:water}{rivier} die verschuift.
\end{solution}

\begin{solution}{exo:g9:maps-of-the-earth:10}
(a) $20\,200 \div 300\,000 \approx 0.067$~s $= 6.7 \times 10^{-2}$~s.
(b) $300\,000 \times 10^{-3} = \qty{300}{km}$. (c) Een piepkleine fout in
de tijd geeft een enorme fout in de afstand: de klokken moeten uiterst
nauwkeurig zijn.
\end{solution}

\begin{solution}{exo:g9:maps-of-the-earth:11}
(a) Afstand tot A: $\sqrt{3.6^2 + 4.8^2} = \sqrt{12.96 + 23.04} =
\sqrt{36} = 6$; afstand tot B: $\sqrt{6.4^2 + 4.8^2} = \sqrt{40.96 +
23.04} = \sqrt{64} = 8$. (b) $(3.6, -4.8)$, het spiegelbeeld ten opzichte
van de lijn AB. (c) Afstand van C tot $(3.6, 4.8)$: $\sqrt{3.6^2 + 5.2^2} =
\sqrt{40} \approx 6.3$; tot $(3.6, -4.8)$: $\sqrt{3.6^2 + 14.8^2} =
\sqrt{232} \approx 15.2$. De ontvanger staat op $(3.6, 4.8)$. (d) Eén
afstand geeft een cirkel (een bol in de ruimte), twee geven twee mogelijke
punten; er is nog een afstand nodig om te kiezen.
\end{solution}

\begin{solution}{exo:g9:maps-of-the-earth:12}
(a) $29 \div 2.1 \approx 14$ keer. (b) $2.1 \times 10.5 \approx
22$~miljoen~\unit{km^2}. (c) Groenland lijkt ongeveer 22~miljoen
\unit{km^2} te beslaan, dicht bij de 29~miljoen van Afrika: op de
mercatorwandkaart lijken de twee ongeveer even groot, hoewel Afrika
14 keer zo groot is. (d) Een \omterm{def:g9:maps-of-the-earth:projection}{oppervlaktegetrouwe projectie}.
\end{solution}

\begin{solution}{pb:g9:maps-of-the-earth:1}
\textbf{1.} $50\,000$~cm $= \qty{500}{m}$.

\textbf{2.} $3.6 \times 500 = \qty{1800}{m} = \qty{1.8}{km}$.

\textbf{3.} B--W: $2.5 \times 500 = \qty{1250}{m}$; W--T: $2.35 \times 500
\approx \qty{1180}{m}$.

\textbf{4.} Naar het noordoosten.

\textbf{5.} \qty{50}{m}.

\textbf{6.} S--B gaat omlaag: $296 - 143 = \qty{153}{m}$.

\textbf{7.} B--W: $335 - 143 = \qty{192}{m}$; W--T: $544 - 335 =
\qty{209}{m}$.

\textbf{8.} $209 \div 1180 \approx 0.18$, ongeveer \qty{18}{\%}.

\textbf{9.} Een rechte lijn van S naar T zou de \omterm{def:g1:landscapes-around-us:water}{rivier} kruisen op een plek
waar geen brug is; het pad gaat omlaag naar de brug bij B.

\textbf{10.} $1.8 + 1.25 + 1.18 \approx \qty{4.2}{km}$ met \qty{4}{km} per
uur: ongeveer 63 minuten; klim $192 + 209 = \qty{401}{m}$: ongeveer
40 minuten extra. In totaal ongeveer 1~uur~43~min.

\textbf{11.} Vanaf T daalt ze af naar B en klimt ze dan van B omhoog naar
S: $296 -
143 = \qty{153}{m}$.

\textbf{12.} Lengte: $1.8 + 1.25 + 1.18 \approx \textbf{4.2 km}$; totale
stijging: $192 + 209 = \textbf{401 m}$.
\end{solution}
